% !TeX program = xelatex
% =====================================================================
%  个人学术简历 -- 张晗（中文）
%
%  编译： xelatex cv-zh.tex  （连续两次以生成页码引用）
%  在 Overleaf 上：菜单 Menu -> Compiler 必须选 XeLaTeX（Overleaf 不读上面
%  那行 magic comment，只认项目设置）。
%
%  内容来源： https://netsyscode.github.io/netsyslab/index.zh.html
%  教育背景依据清华大学网络研究院官网教师页
%  （https://www.insc.tsinghua.edu.cn/info/1157/2458.htm）。
% =====================================================================
\documentclass[11pt]{article}

\newcommand{\cvfooter}{张晗 \textbullet\ 学术简历}

% ctex 必须在 titlesec（由 cvstyle 载入）之前加载，否则章节格式会被覆盖。
% fontset=fandol 用 ctex 自带的开源方正字体，Overleaf 与本机都有，不依赖系统字体。
\usepackage[UTF8,fontset=fandol,scheme=plain]{ctex}
\usepackage{cvstyle}

\setcounter{secnumdepth}{0}
\linespread{1.10}

\begin{document}

% ---------------------------------------------------------------- 页首
\begin{center}
  {\Huge\bfseries\color{accent} 张\ 晗}\quad{\LARGE\bfseries\color{accent} Han Zhang}\\[5pt]
  {\large 副教授}\\[3pt]
  清华大学\ 网络科学与网络空间研究院\\
  中国\ 北京\\[5pt]
  \small
  \href{mailto:zhhan@tsinghua.edu.cn}{zhhan@tsinghua.edu.cn}\ \textbullet\
  \href{https://netsyscode.github.io/netsyslab/}{netsyscode.github.io/netsyslab}\ \textbullet\
  \href{https://scholar.google.com/citations?user=2_wPPU0AAAAJ}{Google Scholar}\ \textbullet\
  \href{https://github.com/netsyscode}{github.com/netsyscode}
\end{center}
\vspace{2pt}

% ------------------------------------------------------------- 研究概述
\section*{研究概述}
\noindent
本人的研究致力于为生产规模的网络系统构建\textbf{测量、验证与态势感知基础设施}，覆盖四类部署场景：
\textbf{（一）互联网}——证书透明度（CT）日志、BGP、RPKI、DNS、Web-PKI；
\textbf{（二）ISP 与运营商网络}——包括 \textbf{FITI / CERNET / CERNET2} 国家级试验设施与运营商骨干网；
\textbf{（三）智算与数据中心基础设施}——超大规模云、Serverless 平台、智能计算中心；
以及\textbf{（四）安全可信机器人}，作为独立但互补的研究方向。

\smallskip
\noindent
研究工作组织为三条主线，共同汇聚为一个统一议程：\emph{下一代互联网、国家级试验设施、
运营智算基础设施以及安全攸关的具身智能系统所需的可观测基底、验证工具链与信任管道}。

\begin{enumerate}[leftmargin=1.8em,label=\textbf{\arabic*.}]
  \item \textbf{网络空间态势感知与测量。} 面向互联网与试验床的可观测性与威胁检测：
  基于 eBPF 与网内遥测的零代码分布式追踪（SIGCOMM'23、SIGCOMM'25、ICNP'25）；
  可编程交换机上的线速加密流量入侵检测（SIGCOMM'25、SIGCOMM'26）；
  基于溯源图的 APT 检测（TIFS'22、TIFS'25）；Web-PKI 与证书误签发监测（WWW'25）；
  互联网拓扑与 Looking-Glass 挖掘（WWW'26）。
  \item \textbf{AI 赋能网络与运维可验证性。} 使运营基础设施可验证、可自诊、可韧性恢复：
  运营商规模的配置验证（TIFS'24、INFOCOM'26）；覆盖网集成验证（CoNEXT'25）；
  路由信任分析（CCS'25、ICNP'26）；测试报文生成（TPDS'23）；
  跨数据中心广域网流量工程（ToN'22、CoNEXT'21）；
  Serverless 钱包耗尽（DoW）攻击与负载验证（CCS'22、ToN'26）；
  面向集合通信的网内前向纠错（INFOCOM'24）。
  \item \textbf{安全可信机器人。} 将同样的测量与验证方法学迁移到具身智能：
  可信多模态感知、智能体与机器人控制器的运行时验证、加密传感器通道异常检测、
  自主车队的信任证书——以对抗鲁棒机器学习（SIGCOMM'25、CCS'21、NDSS'23）、
  溯源图 APT 检测与 PKI 信任为基础。
\end{enumerate}

% ----------------------------------------------------------- 学术计量
\section*{学术计量}
\metricrow{%
  \metric{1{,}730}{引用次数}\hfill
  \metric{17}{h 指数}\hfill
  \metric{29}{i10 指数}\hfill
  \metric{近 100}{论文}\hfill
  \metric{13}{在读学生}}
\vspace{6pt}

\noindent\note{数据来源 Google Scholar，2026 年。已在 ACM SIGCOMM、USENIX NSDI、ACM CCS、NDSS、
ACM CoNEXT、ACM WWW、ACM ICSE、IEEE INFOCOM、IEEE/ACM ToN、IEEE TIFS、IEEE TPDS
等国际顶级会议与期刊发表论文近百篇。}

% -------------------------------------------------------------- 教育背景
\section{教育背景}
\begin{kvlist}
  \kvitem{2013.09 -- 2018.07}{\textbf{工学博士}，计算机科学与技术 \\ 清华大学，中国\ 北京 \quad \note{导师：尹霞 教授}}
  \kvitem{2009.09 -- 2013.07}{\textbf{工学学士}，计算机科学与技术 \\ 吉林大学，中国\ 长春}
\end{kvlist}

% ----------------------------------------------------------- 工作履历
\section{工作履历}
\begin{kvlist}
  \kvitem{现任}{\textbf{副教授}、博士生导师 \\ 清华大学 网络科学与网络空间研究院}
  \kvitem{曾任}{\textbf{助理教授} \\ 清华大学 网络科学与网络空间研究院}
\end{kvlist}

% ---------------------------------------------------------------- 奖励
\section{主要奖励与荣誉}
\begin{itemize}
  \item \textbf{Applied Networking Research Prize（ANRP）}，IRTF / Internet Society，2023。
        \note{国际}
  \item \textbf{中国通信学会科学技术奖一等奖}——《城域云网关键技术及规模应用》，2025。
        \note{国家级学会}
  \item \textbf{NSGA 杰出研究者奖（Best Researcher Award）}，2025。
  \item \textbf{ACM TURC SIGCSE 中国新星奖}，2023。
  \item \textbf{优秀服务奖（Excellent Service Award）}，IEEE/ACM ICA3PP 2023。
\end{itemize}

% --------------------------------------------------- 代表性论文
\section{代表性论文}
\note{\me{加粗}为本人；\corr{} 表示通信作者。完整列表见下一节。}
\begin{publist}

\item \me{Han Zhang}, Xuefeng Liu, Linqiang Qian, Guyue Liu, Tianyu Zhang, Kaiyang
Zhao, Yantu Tong, Zeji Xiao, Dongbiao He, Yahui Li, Ke Ruan, Jilong Wang, Yongqing
Zhu, Xia Yin. ``Towards High-Performance Intrusion Detection with Robustness
Guarantees on Programmable Switches at ISP Scale.'' \venue{ACM SIGCOMM 2026},
pp.~1282--1298.
\insight{首个面向 ISP 的可编程交换机入侵检测体系：数据平面轻量特征提取与
控制平面自适应学习协同设计，在 Tbps 量级线速检测的同时对流量漂移与
对抗性规避保持鲁棒；在真实 ISP 骨干网验证，资源开销相比 CPU/GPU 方案降低数个量级。}

\item Kaiyang Zhao, \me{Han Zhang}\corr, Tianpeng Zhang, Yantu Tong, Zeji Xiao,
Xuefeng Liu, Xia Yin, Jianping Wu. ``Kuiper: Enterprise Cloud Plug-and-Play Network
Management Platform.'' \venue{USENIX NSDI 2027}.
\insight{面向企业云的统一智能运维控制面：在同一套遥测与验证基底上整合
故障管理、安全事件响应与配置管理。}

\item Yantao Geng, \me{Han Zhang}\corr, Zhiheng Wu, Yahui Li, Jilong Wang, Xia Yin.
``Low-Overhead Distributed Application Observation with DeepTrace: Achieving Accurate
Tracing in Production Systems.'' \venue{ACM SIGCOMM 2025}, pp.~1056--1069.
\insight{方法级无侵入追踪，开销仅为侵入式基线的约 $1/10$；
高并发场景下准确率仍超过 95\%；已被数十家企业用于生产环境的故障诊断与资源优化。}

\item \me{Han Zhang}, Guyue Liu\corr, Xingang Shi, Yahui Li, Jilong Wang, Yongqing
Zhu, Ke Ruan, Jie Liang, Xia Yin. ``Achieving High-Speed and Robust Encrypted Traffic
Anomaly Detection with Programmable Switches.'' \venue{ACM SIGCOMM 2025},
pp.~1254--1256.
\insight{双平面特征提取 + 轻量单类分类模型，化解深度学习精度与交换芯片
资源约束之间的长期矛盾；已在生产环境运行两年，发现 10 余起重大攻击事件，
为两家 ISP 保护超过 1000 万台服务器。}

\item Junxian Shen, \me{Han Zhang}\corr, Yang Xiang, Xingang Shi, Xinrui Li, Yunxi
Shen, Zijian Zhang, Sicheng Wang, Yanhui Yang, Yu Zhou, Mingwei Xu, Jilong Wang.
``Network-Centric Distributed Tracing with DeepFlow: Troubleshooting Your
Microservices in Zero Code.'' \venue{ACM SIGCOMM 2023}, pp.~420--437.
\insight{以网络为中心的追踪平面与隐式上下文传播，完全无需修改应用代码。
已开源并成为云原生可观测性的事实标准（CNCF Sandbox，4000+ stars，部署于数万台主机）。}

\item Yunze Wei, Xingang Shi, Minghao Zhang, Zhaozhen Wang, Qiang Ma, \me{Han Zhang},
Yahui Li, Xia Yin. ``Parade: Unified and Data-Driven Prefix-AS Semantics for BGP
Anomaly Detection.'' \venue{ACM CCS 2026}.
\insight{统一的数据驱动前缀—AS 语义模型，把碎片化的 BGP 观测转化为
可用于异常检测的一致性基准事实。}

\item Junxian Shen, \me{Han Zhang}\corr, Yantao Geng, Jiawei Li, Jilong Wang, Mingwei
Xu. ``Gringotts: Fast and Accurate Internal Denial-of-Wallet Detection for Serverless
Computing.'' \venue{ACM CCS 2022}, pp.~2627--2641.
\insight{首次系统性研究 Serverless 平台上的内部钱包耗尽（DoW）攻击；
性能开销低于 1.1\%，平均检测时延 1.86 秒，准确率超过 95.75\%。}

\item Dongqi Han, Zhiliang Wang\corr, Wenqi Chen, Ying Zhong, Su Wang, \me{Han Zhang},
Jiahai Yang, Xingang Shi, Xia Yin. ``DeepAID: Interpreting and Improving Deep
Learning-based Anomaly Detection in Security Applications.'' \venue{ACM CCS 2021},
pp.~3197--3217.
\insight{面向安全领域\emph{无监督}深度异常检测模型的解释框架，
使运维人员能够诊断模型误判并显著降低误报率。}

\item Dongqi Han, Zhiliang Wang\corr, Wenqi Chen, Kai Wang, Rui Yu, Su Wang, \me{Han
Zhang}, et al. ``Anomaly Detection in the Open World: Normality Shift Detection,
Explanation, and Adaptation.'' \venue{NDSS 2023}.
\insight{将入侵检测从离线分析推进到正常态漂移条件下的在线、ISP 规模运行。}

\item Yantao Geng, \me{Han Zhang}\corr, Zhiheng Wu, Yahui Li, Jilong Wang, Xia Yin.
``Non-Intrusive Distributed Tracing with Method-Level Delay Estimation for
Microservices Troubleshooting.'' \venue{ACM/IEEE ICSE 2026}, pp.~1661--1673.
\insight{从网络侧观测结合轻量运行时线索推断方法级时延，
故障定位精度与侵入式方案持平，而工程代价大幅降低。}

\item Tianyu Zhang, \me{Han Zhang}\corr, Yunze Wei, Yahui Li, Xingang Shi, Jilong
Wang, Xia Yin. ``ACME++: A Secure Authorization Mechanism for ACME Clients in the Web PKI
Ecosystem.''
\venue{ACM WWW 2025}, pp.~1058--1067.
\insight{披露 \emph{ACME Authz 缓存攻击}——攻击者无需控制域名即可获取
欺诈证书（在 Let's Encrypt 上实证），并提出开销可忽略的账户绑定缓解机制。}

\item Yunze Wei, Xingang Shi, \me{Han Zhang}\corr, Tianyu Zhang, Yahui Li, Xia Yin.
``GlassMiner: Mining Looking Glass Services via Structure-Semantics Fusion for Web
Observability.'' \venue{ACM WWW 2026}, pp.~7576--7587.
\insight{联合建模 Looking Glass 门户的 DOM 结构与自然语言语义，
在数千个运营商门户上实现洲际规模的 BGP 路径可观测性。}

\item XinZhe Liu, Yahui Li\corr, \me{Han Zhang}, Xia Yin, Xingang Shi, Zhiliang Wang,
Gang Ren, Jilong Wang, Jiangyuan Yao. ``Scalable and Interpretable Overlay Network
Checking via Ensemble Verification.'' \venue{ACM CoNEXT 2025}.
\insight{首个面向 overlay/underlay 网络的控制平面验证器；
部署于全国性 FITI 试验设施，发现 22 处配置错误，相比 Batfish+ 最高提速 $102\times$。}

\item \me{Han Zhang}, Xingang Shi\corr, Xia Yin, Jilong Wang, Zhiliang Wang, Yingya
Guo, Tian Lan. ``Boosting Bandwidth Availability over Inter-DC WAN.''
\venue{ACM CoNEXT 2021}, pp.~297--312.
\insight{概率化带宽可用性流量工程：在线准入控制提速 $30\times$，
相比 FFC 与 TEAVAR 多满足 23\%--60\% 的带宽可用性需求；
在腾讯与中国电信真实拓扑上验证。}

\item Kaiyang Zhao, \me{Han Zhang}\corr, Yao Tong, Yahui Li, Xingang Shi, Zhiliang
Wang, Xia Yin, Jianping Wu. ``SRmesh: Deterministic and Efficient Diagnosis of Latency
Bottleneck Links in SRv6 Networks.'' \venue{IEEE ICNP 2025}, pp.~1--12.
\insight{将确定性带内遥测嵌入分段路由报头，在探测开销有界的前提下
实现逐链路时延归因；在多厂商 SRv6 试验床上召回率 100\%。}

\item Yahui Li, \me{Han Zhang}\corr, Jilong Wang, Xingang Shi, Xia Yin, Zhiliang Wang,
Jiankun Hu, Congcong Miao, Jianping Wu. ``Proactively Verifying Quantitative Network
Policy across Unsafe and Unreliable Environments.'' \venue{IEEE Transactions on
Information Forensics and Security}, vol.~19, pp.~10099--10113, 2024.
\insight{把形式化验证引入此前依赖人工调参的定量网络策略，
扩展到运营商规模的状态空间，同时保持反例的可解释性。}

\end{publist}

% ------------------------------------------------------ 论文全表
\section{论文列表（按年份）}
\note{\me{加粗}为本人；\corr{} 表示通信作者。}

\subsection*{2027}
\begin{publist}
\item Kaiyang Zhao, \me{Han Zhang}\corr, Tianpeng Zhang, Yantu Tong, Zeji Xiao,
Xuefeng Liu, Xia Yin, Jianping Wu. ``Kuiper: Enterprise Cloud Plug-and-Play Network
Management Platform.'' \venue{USENIX NSDI 2027}. \note{（待刊）}
\end{publist}

\subsection*{2026}
\begin{publist}
\item \me{Han Zhang}, Xuefeng Liu, Linqiang Qian, Guyue Liu, Tianyu Zhang, Kaiyang
Zhao, Yantu Tong, Zeji Xiao, Dongbiao He, Yahui Li, Ke Ruan, Jilong Wang, Yongqing
Zhu, Xia Yin. ``Towards High-Performance Intrusion Detection with Robustness
Guarantees on Programmable Switches at ISP Scale.'' \venue{ACM SIGCOMM 2026},
pp.~1282--1298.
\item Yunze Wei, Xingang Shi, Minghao Zhang, Zhaozhen Wang, Qiang Ma, \me{Han Zhang},
Yahui Li, Xia Yin. ``Parade: Unified and Data-Driven Prefix-AS Semantics for BGP
Anomaly Detection.'' \venue{ACM CCS 2026}. \note{（待刊）}
\item Yunze Wei, Xingang Shi, \me{Han Zhang}\corr, Tianyu Zhang, Yahui Li, Xia Yin.
``GlassMiner: Mining Looking Glass Services via Structure-Semantics Fusion for Web
Observability.'' \venue{ACM WWW 2026}, pp.~7576--7587.
\item Yantao Geng, \me{Han Zhang}\corr, Zhiheng Wu, Yahui Li, Jilong Wang, Xia Yin.
``Non-Intrusive Distributed Tracing with Method-Level Delay Estimation for
Microservices Troubleshooting.'' \venue{ACM/IEEE ICSE 2026}, pp.~1661--1673.
\item Junxian Shen, \me{Han Zhang}\corr, Wenxue Lin, Yantao Geng, Jiawei Li, Jilong
Wang, Mingwei Xu. ``Clover: Workload Verification for Real-Time Detection of
Contention-Induced Slowdowns in Serverless Platforms.'' \venue{IEEE/ACM Transactions
on Networking}, vol.~34, pp.~4700--4715, 2026.
\item Kaiyang Zhao, \me{Han Zhang}\corr, Yao Tong, Yahui Li, Xingang Shi, Zhiliang
Wang, Xia Yin, Jianping Wu. ``Efficient Slice-Parallel Distributed Probing in SRv6
Networks.'' \venue{IEEE/ACM Transactions on Networking}, 2026.
\item XinZhe Liu, Yahui Li, \me{Han Zhang}, Renrui Tian, Xia Yin, Xingang Shi,
Zhiliang Wang, Gang Ren, Jilong Wang, Jiangyuan Yao. ``Comprehensive Network
Configuration Verification via Effective Environment Reduction.''
\venue{IEEE INFOCOM 2026}, pp.~1--10.
\item Letong Sun, Xingang Shi, Fengyan Han, Liyuan Chang, Xia Yin, Zhiliang Wang,
\me{Han Zhang}. ``A Comparative Approach for Studying Control-Plane and Data-Plane
Mismatch in Interdomain Routing.'' \venue{IEEE ICNP 2026}. \note{（待刊）}
\item Letong Sun, Xingang Shi, Fengyan Han, Xia Yin, Zhiliang Wang, \me{Han Zhang}.
``The Anatomy of ROV Rollout: A Sub-AS and Longitudinal Perspective.''
\venue{IEEE ICNP 2026}. \note{（待刊）}
\item Weihan Chen, Zhiliang Wang, \me{Han Zhang}, Xia Yin, Xingang Shi.
``Delay-Aware Service Function Chain Embedding and Scheduling in Edge Clouds.''
\venue{IEEE Transactions on Parallel and Distributed Systems}, vol.~37, no.~11,
pp.~2539--2553, 2026.
\item Shuhuan Wang, Zhiyu Wang, Hao Sun, Zhi Jiang, Xia Yin, Dongbin Han, \me{Han
Zhang}, Xingang Shi, Jiahai Yang. ``The Evolution of Provenance-Based Intrusion
Detection Systems: A Review of Approaches, Limitations, and Future Directions.''
\venue{IEEE Network}, 2026.
\item 李秋香, \me{张晗}\corr, 曲豫宾, 吴建平. ``MASTER：基于知识注入与意图强化的大语言模型
漏洞验证代码生成框架.'' \venue{电子学报}, vol.~54, no.~4, pp.~1802--1822, 2026.
\end{publist}

\subsection*{2025}
\begin{publist}
\item Tianyu Zhang, \me{Han Zhang}\corr, Yunze Wei, Yahui Li, Xingang Shi, Jilong
Wang, Xia Yin. ``ACME++: A Secure Authorization Mechanism for ACME Clients in the Web PKI
Ecosystem.''
\venue{ACM WWW 2025}, pp.~1058--1067.
\item Yantao Geng, \me{Han Zhang}\corr, Zhiheng Wu, Yahui Li, Jilong Wang, Xia Yin.
``Low-Overhead Distributed Application Observation with DeepTrace: Achieving Accurate
Tracing in Production Systems.'' \venue{ACM SIGCOMM 2025}, pp.~1056--1069.
\item \me{Han Zhang}, Guyue Liu\corr, Xingang Shi, Yahui Li, Jilong Wang, Yongqing
Zhu, Ke Ruan, Jie Liang, Xia Yin. ``Achieving High-Speed and Robust Encrypted Traffic
Anomaly Detection with Programmable Switches.'' \venue{ACM SIGCOMM 2025},
pp.~1254--1256.
\item XinZhe Liu, Yahui Li\corr, \me{Han Zhang}, Xia Yin, Xingang Shi, Zhiliang Wang,
Gang Ren, Jilong Wang, Jiangyuan Yao. ``Scalable and Interpretable Overlay Network
Checking via Ensemble Verification.'' \venue{ACM CoNEXT 2025}.
\item Renrui Tian, Yahui Li\corr, \me{Han Zhang}, XinZhe Liu, Xia Yin, Xingang Shi,
Zhiliang Wang, Jilong Wang. ``Poster: ERIS --- Evaluating ROV via ICMPv6 Rate Limiting
Side Channels.'' \venue{ACM CCS 2025}.
\item Kaiyang Zhao, \me{Han Zhang}\corr, Yao Tong, Yahui Li, Xingang Shi, Zhiliang
Wang, Xia Yin, Jianping Wu. ``SRmesh: Deterministic and Efficient Diagnosis of Latency
Bottleneck Links in SRv6 Networks.'' \venue{IEEE ICNP 2025}, pp.~1--12.
\item Zhaozhen Wang, Xingang Shi\corr, Haijun Geng, \me{Han Zhang}, Xia Yin, Zhiliang
Wang. ``On Non-Commutative Routing.'' \venue{IEEE INFOCOM 2025}, pp.~1--10.
\item Zhiyun Tang, Yahui Li, Xingang Shi, Jilong Wang, Yongqing Zhu, Ke Ruan, Xia Yin,
\me{Han Zhang}\corr. ``Probabilistic Analysis of Overload-Free Property for Critical
Traffic.'' \venue{IEEE/ACM IWQoS 2025}, pp.~1--6.
\item Su Wang, Hongbin Sun, Zhiliang Wang, Tao Zhou, Xia Yin, Dongqi Han, \me{Han
Zhang}, Xingang Shi, Jiahai Yang. ``End-to-End Attack Scene Reconstruction in a Host
with Rules and Anomaly-Based Detection Models.'' \venue{IEEE Transactions on
Information Forensics and Security}, 2025.
\item Ying Tian, Zhiliang Wang\corr, Xia Yin, Xingang Shi, Jiahai Yang, \me{Han
Zhang}\corr. ``Centralized Network Utility Maximization with Accelerated Gradient
Method.'' \venue{IEEE Transactions on Networking}, 2025.
\item Xingang Shi, Zitong Jin\corr, Bin Xiong, Xinyao Huang, Xiaotian Xi, Dan Li,
\me{Han Zhang}, Zhiliang Wang, Xia Yin. ``HELA: Inferring AS Relationships with a
Hybrid of Empirical and Learning Algorithms.'' \venue{IEEE Transactions on
Networking}, 2025.
\item Weiting Zhang, Jiadong Ren, Tao Zheng, Ruibin Guo, \me{Han Zhang}, Shiwen Mao,
Hongke Zhang\corr. ``GenNet: Computing-Efficient Generative AI for Deterministic
Transmission Scheduling in 6G Networks.'' \venue{IEEE Communications Magazine},
vol.~63, no.~8, pp.~48--55, 2025.
\item Yanxu Fu, Pei Zhang\corr, \me{Han Zhang}, Xiaohong Huang, Kun Xie, Dandan Li.
``BGPFlow: Flow-Based Feature Extraction for BGP Anomaly Detection.''
\venue{ACM SIGCOMM NGNO Workshop 2025}, pp.~13--20.
\end{publist}

\subsection*{2024}
\begin{publist}
\item Yi Qiao, \me{Han Zhang}\corr, Jilong Wang. ``NetFEC: In-Network FEC Encoding
Acceleration for Latency-Sensitive Multimedia Applications.''
\venue{IEEE INFOCOM 2024}, pp.~2348--2357.
\item Yahui Li, \me{Han Zhang}\corr, Jilong Wang, Xingang Shi, Xia Yin, Zhiliang Wang,
Jiankun Hu, Congcong Miao, Jianping Wu. ``Proactively Verifying Quantitative Network
Policy across Unsafe and Unreliable Environments.'' \venue{IEEE Transactions on
Information Forensics and Security}, vol.~19, pp.~10099--10113, 2024.
\item XinZhe Liu, Yahui Li\corr, \me{Han Zhang}, Xia Yin, Xingang Shi, Zhiliang Wang,
Gang Ren, Jiangyuan Yao. ``Poster: Scalable and Interpretable Multilayer Overlay
Network Checking via Ensemble Verification.'' \venue{ACM SIGCOMM 2024 Posters \&
Demos}.
\item Weihan Chen, Zhiliang Wang\corr, \me{Han Zhang}, Xia Yin, Xingang Shi.
``Cost-Efficient Flow Migration for SFC Dynamical Scheduling in Geo-Distributed
Clouds.'' \venue{Computer Networks}, vol.~249, 2024.
\item Jun He, Yahui Li\corr, \me{Han Zhang}\corr, Changqing An, Jilong Wang. ``ROV-GD:
Improving the Measurement of ROV Deployment Using Graph Difference.''
\venue{IEEE ISCC 2024}, pp.~1--7.
\item Kunling He, Fenghua Li, Jessie Hui Wang, \me{Han Zhang}, Yiren Zhao. ``AB-TCAD:
An Access Behavior-Based Two-Stage Compromised Account Detection Framework.''
\venue{IFIP Networking 2024}, pp.~104--112.
\end{publist}

\subsection*{2023}
\begin{publist}
\item Junxian Shen, \me{Han Zhang}\corr, Yang Xiang, Xingang Shi, Xinrui Li, Yunxi
Shen, Zijian Zhang, et al. ``Network-Centric Distributed Tracing with DeepFlow:
Troubleshooting Your Microservices in Zero Code.'' \venue{ACM SIGCOMM 2023},
pp.~420--437.
\item Dongqi Han, Zhiliang Wang\corr, Wenqi Chen, Kai Wang, Rui Yu, Su Wang, \me{Han
Zhang}, et al. ``Anomaly Detection in the Open World: Normality Shift Detection,
Explanation, and Adaptation.'' \venue{NDSS 2023}.
\item Heng Yu, \me{Han Zhang}\corr, Junxian Shen, Yantao Geng, et al. ``Serpens: A
High-Performance FaaS Platform for Network Functions.'' \venue{IEEE Transactions on
Parallel and Distributed Systems}, vol.~34, no.~8, pp.~2448--2463, 2023.
\item Yahui Li, \me{Han Zhang}\corr, Jilong Wang, Xia Yin, Xingang Shi, Jianping Wu.
``A General Approach to Generate Test Packets with Network Configurations.''
\venue{IEEE Transactions on Parallel and Distributed Systems}, vol.~34, no.~4,
pp.~1362--1375, 2023.
\item Pei Zhang, Fangzhou He, \me{Han Zhang}\corr, Jiankun Hu, et al. ``Real-Time
Malicious Traffic Detection with Online Isolation Forest over SD-WAN.''
\venue{IEEE Transactions on Information Forensics and Security}, vol.~18,
pp.~2076--2090, 2023.
\item Tianbo Wang, Huacheng Li, Chunhe Xia\corr, \me{Han Zhang}, Pei Zhang. ``From the
Dialectical Perspective: Modeling and Exploiting of Hybrid Worm Propagation.''
\venue{IEEE Transactions on Information Forensics and Security}, vol.~18,
pp.~1610--1624, 2023.
\item Yantao Geng, \me{Han Zhang}\corr, Jilong Wang. ``Delay-Based Congestion Control
for Cross-Datacenter Networks.'' \venue{IEEE/ACM IWQoS 2023}.
\item Jiawei Li, \me{Han Zhang}\corr, Jilong Wang. ``PUVAR: Minimize Idle Resource SLO
Violations by Uncertainty-Aware Scheduling in Cloud Platforms.''
\venue{IEEE ICWS 2023}.
\item Ming Wang, Yahui Li, \me{Han Zhang}, Anlun Hong, Jun He, Jilong Wang. ``Knocking
Cells: Latency-Based Identification of IPv6 Cellular Addresses on the Internet.''
\venue{IEEE ISCC 2023}.
\item Allen Hong, Yahui Li\corr, \me{Han Zhang}\corr, Ming Wang, Changqing An, Jilong
Wang. ``A Cheap and Accurate Delay-Based IP Geolocation Method Using Machine Learning
and Looking Glass.'' \venue{IFIP Networking 2023}.
\item Jun He, Yahui Li\corr, \me{Han Zhang}\corr, Ming Wang, Changqing An, Jilong
Wang. ``Be Careful of Your Neighbors: Injected Sub-Prefix Hijacking Invisible to
Public Monitors.'' \venue{IEEE ICC 2023}.
\item Shuanghong Yu, Fenghua Li\corr, \me{Han Zhang}, Kunling He, Xingkun Yao.
``GapReplay: A High-Accuracy Packet Replayer.'' \venue{IEEE ICC 2023}.
\end{publist}

\subsection*{2022}
\begin{publist}
\item Junxian Shen, \me{Han Zhang}\corr, Yantao Geng, Jiawei Li, Jilong Wang, Mingwei
Xu. ``Gringotts: Fast and Accurate Internal Denial-of-Wallet Detection for Serverless
Computing.'' \venue{ACM CCS 2022}, pp.~2627--2641.
\item \me{Han Zhang}, Xia Yin, Xingang Shi\corr, Jilong Wang, Zhiliang Wang, Yingya
Guo, Tian Lan, Yahui Li, Yongqing Zhu, Ke Ruan, Haijun Geng. ``Achieving High
Availability in Inter-DC WAN Traffic Engineering.'' \venue{IEEE/ACM Transactions on
Networking}, 2022.
\item Su Wang, Zhiliang Wang\corr, Tao Zhou, Hongbin Sun, Xia Yin, Dongqi Han,
\me{Han Zhang}, et al. ``THREATRACE: Detecting and Tracing Host-Based Threats in Node
Level through Provenance Graph Learning.'' \venue{IEEE Transactions on Information
Forensics and Security}, vol.~17, pp.~3972--3987, 2022.
\item Yi Qiao, Menghao Zhang, Yu Zhou, Xiao Kong, \me{Han Zhang}\corr, et al. ``NetEC:
Accelerating Erasure Coding Reconstruction with In-Network Aggregation.''
\venue{IEEE Transactions on Parallel and Distributed Systems}, vol.~33, no.~10,
pp.~2571--2583, 2022.
\item Zongyi Zhao, Xingang Shi\corr, Zhiliang Wang, Qi Li, \me{Han Zhang}, Xia Yin.
``Efficient and Accurate Flow Record Collection with HashFlow.'' \venue{IEEE
Transactions on Parallel and Distributed Systems}, vol.~33, no.~5, pp.~1069--1083,
2022.
\item Ying Tian, Zhiliang Wang\corr, Xia Yin, Xingang Shi, Jiahai Yang, \me{Han
Zhang}. ``Centralized Network Utility Maximization with Accelerated Gradient Method.''
\venue{IEEE ICNP 2022}.
\item Heng Yu, Junxian Shen, \me{Han Zhang}\corr, Jilong Wang, Congcong Miao, Mingwei
Xu. ``Scorpius: Proactive Code Preparation to Accelerate Function Startup.''
\venue{IEEE/ACM IWQoS 2022}.
\item Shangbin Han, Qianhong Wu, \me{Han Zhang}\corr, Bo Qin. ``Light-Weight
Unsupervised Anomaly Detection for Encrypted Malware Traffic.'' \venue{IEEE DSC 2022}.
\item Shangbin Han, Qianhong Wu, \me{Han Zhang}, Bo Qin, Jiangyuan Yao, Willy Susilo.
``Autoperman: Automatic Network Traffic Anomaly Detection with Ensemble Learning.''
\venue{ICAIS 2022}.
\item Yongkang Chen, Ming Zhang\corr, Jin Li, Xiaohui Kuang, Xurong Zhang, \me{Han
Zhang}. ``Dynamic and Diverse Transformations for Defending against Adversarial
Examples.'' \venue{IEEE TrustCom 2022}.
\end{publist}

\subsection*{2021}
\begin{publist}
\item \me{Han Zhang}, Xingang Shi\corr, Xia Yin, Jilong Wang, Zhiliang Wang, Yingya
Guo, Tian Lan. ``Boosting Bandwidth Availability over Inter-DC WAN.''
\venue{ACM CoNEXT 2021}, pp.~297--312.
\item Dongqi Han, Zhiliang Wang\corr, Wenqi Chen, Ying Zhong, Su Wang, \me{Han Zhang},
Jiahai Yang, Xingang Shi, Xia Yin. ``DeepAID: Interpreting and Improving Deep
Learning-Based Anomaly Detection in Security Applications.'' \venue{ACM CCS 2021},
pp.~3197--3217.
\item Shangbin Han, Qianhong Wu, \me{Han Zhang}\corr, Bo Qin, Jiankun Hu, et al.
``Log-Based Anomaly Detection with Robust Feature Extraction and Online Learning.''
\venue{IEEE Transactions on Information Forensics and Security}, vol.~16,
pp.~2300--2311, 2021.
\item Yingya Guo, Weipeng Wang, \me{Han Zhang}, Wenzhong Guo, Zhiliang Wang, Ying
Tian, et al. ``Traffic Engineering in Hybrid Software-Defined Network via
Reinforcement Learning.'' \venue{Journal of Network and Computer Applications},
vol.~189, 2021.
\item Weihan Chen, Zhiliang Wang\corr, \me{Han Zhang}, Xia Yin, Xingang Shi.
``Cost-Efficient Dynamic Service Function Chain Embedding in Edge Clouds.''
\venue{IEEE CNSM 2021}.
\item Ying Tian, Zhiliang Wang\corr, Xia Yin, Xingang Shi, Jiahai Yang, \me{Han
Zhang}, Yingya Guo, et al. ``Traffic Engineering with Segment Routing Considering
Probabilistic Failures.'' \venue{IEEE CNSM 2021}.
\item Haijun Geng, \me{Han Zhang}\corr, Yangyang Zhang. ``Efficient Routing Protection
Algorithm in Large-Scale Networks.'' \venue{Computers, Materials \& Continua},
vol.~66, no.~2, pp.~1733--1744, 2021.
\end{publist}

\subsection*{2020}
\begin{publist}
\item Yahui Li, Zhiliang Wang\corr, Xia Yin, Xingang Shi, Jianping Wu, Fangdan Ye,
Jiangyuan Yao, \me{Han Zhang}. ``Assisting Reachability Verification of Network
Configurations Updates with NUV.'' \venue{Computer Networks}, vol.~177, 2020.
\item Haijun Geng, \me{Han Zhang}, Xingang Shi\corr, Zhiliang Wang, Xia Yin.
``Efficient Computation of Loop-Free Alternates.'' \venue{Journal of Network and
Computer Applications}, vol.~151, 2020.
\item Haijun Geng, \me{Han Zhang}, Xingang Shi\corr, et al. ``A Hybrid Link Protection
Scheme for Ensuring Network Service Availability in Link-State Routing Networks.''
\venue{Journal of Communications and Networks}, vol.~22, no.~1, pp.~46--60, 2020.
\item Weihan Chen, Zhiliang Wang, \me{Han Zhang}, Xia Yin, Xingang Shi, Lei Sun.
``Joint Optimization of Service Function Chain Elastic Scaling and Routing.''
\venue{IEEE INFOCOM Workshops 2020}.
\end{publist}

\subsection*{2019}
\begin{publist}
\item \me{Han Zhang}, Haijun Geng, Yahui Li, et al. ``DA\&FD: Deadline-Aware and Flow
Duration-Based Rate Control for Mixed Flows in DCNs.'' \venue{IEEE/ACM Transactions on
Networking}, vol.~27, no.~6, pp.~2458--2471, 2019.
\item Zhiliang Wang, \me{Han Zhang}\corr, Xingang Shi, et al. ``Efficient Scheduling of
Weighted Coflows in Data Centers.'' \venue{IEEE Transactions on Parallel and
Distributed Systems}, vol.~30, no.~9, pp.~2003--2017, 2019.
\item Yahui Li, Xia Yin, Zhiliang Wang\corr, Jiangyuan Yao, Xingang Shi, Jianping Wu,
\me{Han Zhang}, et al. ``A Survey on Network Verification and Testing with Formal
Methods: Approaches and Challenges.'' \venue{IEEE Communications Surveys \&
Tutorials}, vol.~21, no.~1, pp.~940--969, 2019.
\item Yahui Li, Zhiliang Wang\corr, Jiangyuan Yao, Xia Yin, Xingang Shi, Jianping Wu,
\me{Han Zhang}. ``MSAID: Automated Detection of Interference in Multiple SDN
Applications.'' \venue{Computer Networks}, vol.~153, 2019.
\item Yingya Guo, Zhiliang Wang\corr, \me{Han Zhang}, Xia Yin, Xingang Shi, Jianping
Wu. ``Joint Optimization of Tasks Placement and Routing to Minimize Coflow Completion
Time.'' \venue{Journal of Network and Computer Applications}, vol.~135, 2019.
\end{publist}

\subsection*{2018 年及以前}
\begin{publist}
\item Xiao Ma, Chuang Lin, \me{Han Zhang}\corr, Jianwei Liu. ``Energy-Aware
Computation Offloading of IoT Sensors in Cloudlet-Based Mobile Edge Computing.''
\venue{Sensors}, 2018.
\item Zhifeng Liu, Xia Yin, Zhiliang Wang, Xingang Shi, Yiqiang Hua, Xiaohong Huang,
Liang Wei, Jiangyuan Yao, \me{Han Zhang}, et al. ``Design and Construction of
Hierarchical and Cross-Regional SDN Verification and Demonstration System.''
\venue{Telecommunications Science}, 2018.
\item \me{Han Zhang}, Xingang Shi\corr, Xia Yin, Zhiliang Wang. ``Yosemite: Efficient
Scheduling of Weighted Coflows in Data Centers.'' \venue{IEEE ICNP 2017}.
\item \me{Han Zhang}, Xingang Shi\corr, Yingya Guo, Zhiliang Wang, Xia Yin. ``More
Load, More Differentiation --- Let More Flows Finish before Deadline in Data Center
Networks.'' \venue{Computer Networks}, vol.~127, 2017.
\item \me{Han Zhang}, Xingang Shi, Yingya Guo, Haijun Geng, Zhiliang Wang, Xia Yin.
``Joint Source Selection and Transfer Optimization for Erasure Coding Storage
System.'' \venue{IEEE IPCCC 2017}.
\item \me{Han Zhang}, Xingang Shi, Xia Yin, Zhiliang Wang, Yingya Guo. ``FDRC: Flow
Duration Time Based Rate Control in Data Center Networks.'' \venue{IEEE/ACM IWQoS
2016}.
\item M. Gharaibeh, \me{Han Zhang}, C. Papadopoulos, J. Heidemann. ``Assessing
Co-Locality of IP Blocks.'' \venue{IEEE INFOCOM Workshops 2016}.
\item Zongyi Zhao, Qi Li, Mingwei Xu, Xingang Shi, \me{Han Zhang}. ``Reduce Completion
Time and Guarantee Throughput by Transport with Slight Congestion.''
\venue{IEEE ICC 2016}.
\item \me{Han Zhang}, Xingang Shi, Xia Yin, Fengyuan Ren, Zhiliang Wang. ``More Load,
More Differentiation --- A Design Principle for Deadline-Aware Congestion Control.''
\venue{IEEE INFOCOM 2015}.
\item Haijun Geng, Xingang Shi, Xia Yin, Zhiliang Wang, \me{Han Zhang}. ``Algebra and
Algorithms for Efficient and Correct Multipath QoS Routing in Link State Networks.''
\venue{IEEE/ACM IWQoS 2015}.
\item Haijun Geng, Xingang Shi, Xia Yin, Zhiliang Wang, \me{Han Zhang}. ``An Efficient
Link Protection Scheme for Link-State Routing Networks.'' \venue{IEEE ICC 2015}.
\item Yingya Guo, Zhiliang Wang, Xia Yin, Xingang Shi, Jianping Wu, \me{Han Zhang}.
``Incremental Deployment for Traffic Engineering in Hybrid SDN Network.''
\venue{IEEE IPCCC 2015}.
\item \me{Han Zhang}, C. Papadopoulos. ``BotTalker: Generating Encrypted, Customizable
C\&C Traces.'' \venue{IEEE HST 2015}.
\item Haijun Geng, Xingang Shi, Xia Yin, Zhiliang Wang, \me{Han Zhang}, Jiangyuan Yao.
``A Hybrid Link Protection Scheme for Link-State Routing Networks.''
\venue{IEEE IPCCC 2014}.
\item Haijun Geng, Xingang Shi, Xia Yin, Zhiliang Wang, \me{Han Zhang}, Jiangyuan Yao.
``Let More Nodes Have a Second Choice.'' \venue{IEEE IPCCC 2014}.
\end{publist}

% --------------------------------------------------------------- 科研项目
\section{科研项目}
\begin{itemize}
  \item \textbf{国家自然科学基金面上项目}（No.~62572269）——
        基于网络行为分析的分布式微服务架构应用异常检测技术研究。\textbf{项目负责人}。
  \item \textbf{国家自然科学基金青年项目}（No.~62002009）——
        面向批量传输的数据中心广域网资源分配研究。\textbf{项目负责人}。
  \item \textbf{国家自然科学基金重大项目}（No.~62394322）——
        智算融合网络组网与路由关键技术研究。\textbf{子课题负责人}。
  \item \textbf{国家自然科学基金重点项目}（No.~61932011）——
        基于区块链的物联网安全技术研究。\textbf{子课题负责人}。
  \item \textbf{工信部重点研发项目}（No.~TC220A04A-78）——
        高精度北斗车规级定位芯片及终端定位系统。\textbf{项目负责人}。
  \item \textbf{科技部国家重点研发计划}（No.~2018YFB1800400）——
        大规模安全可信寻址路由关键技术。参与。
  \item \textbf{企业合作项目}（均为\textbf{项目负责人}）：中国电信——大规模广域网感知与优化技术研究；
        腾讯科技——大规模微服务系统可观测性技术；深信服——应用行为感知的广域网传输优化技术；
        绿盟科技——安全试验设施关键技术研究。
\end{itemize}

% ---------------------------------------------------- 系统与开源
\section{生产系统与开源成果}
\begin{itemize}
  \item \textbf{DeepFlow}——面向云原生与 AI 智能体流水线的零代码 eBPF 可观测性平台。
        \textbf{CNCF Sandbox 项目}；GitHub 4000+ stars、450+ forks；部署于数万台主机。
        \note{SIGCOMM'23}
  \item \textbf{DeepTrace}——面向智算中心与大规模微服务的无侵入分布式应用追踪；
        同一套观测流水线延伸至机器人 / 具身智能运行时安全。170+ stars。
        \note{SIGCOMM'25、ICSE'26}
  \item \textbf{SiteGuard / 智算中心出入口防护}——智算中心出入口流量监测与机器学习威胁检测，
        可运行于可编程交换机（P4/Tofino）、DPU、SmartNIC 与 eBPF，达到 Tbps 量级线速。
        \note{SIGCOMM'25、SIGCOMM'26}
  \item \textbf{Kuiper / 智算中心智能运维}——整合故障管理、安全事件响应、配置管理与网络验证的
        全栈 AIOps 平台；在生产 GPU 集群环境中将 MTTR 降低 $5\times$。
        \note{NSDI'27、ICNP'25、ToN'26}
  \item \textbf{NetVerify 套件}——面向配置、覆盖网与定量策略的静态验证与仿真工具族，
        已应用于 FITI 与 CERNET。\note{TIFS'24、CoNEXT'25、INFOCOM'26}
  \item \textbf{SRmesh 探测工具包}——切片并行 SRv6 探测，以 $O(1)$ 轮次获得全网状时延图；
        在国家级 FITI 试验设施上验证。\note{ICNP'25、ToN'26}
  \item \textbf{IScan}——持续的互联网主动 / 被动探测（IP 空间、BGP/RPKI 路由、CT 日志与 Web-PKI），
        汇聚为可查询的互联网知识库。\note{WWW'25、WWW'26}
  \item \textbf{STRobot}——面向自主机器人与具身智能体的安全优先运行时观测与防护框架。
        \note{ICSE'26}
  \item \textbf{公开数据集}——GlassMiner Looking-Glass 语料（WWW'26）；
        带标注的加密流量 IDS 基准数据集（SIGCOMM'25/'26、CCS'21、NDSS'23）。
\end{itemize}

% -------------------------------------------------------------- 成果转化
\section{成果转化与国家级影响}
\begin{itemize}
  \item \textbf{中国电信骨干网}——线速加密流量异常检测；大规模广域网感知与优化。
  \item \textbf{未来互联网试验设施（FITI）}——可编程交换机入侵检测与分段路由诊断在
        国家下一代互联网试验设施上部署；配置与覆盖网验证。
  \item \textbf{CERNET / CERNET2}——基于 DeepFlow 的可观测性与路由异常分析，
        覆盖中国教育和科研计算机网。
  \item \textbf{腾讯云}——大规模微服务系统可观测性；跨数据中心广域网流量工程联合设计。
  \item \textbf{深信服}——面向企业 SD-WAN 产品的应用感知广域网传输优化。
  \item \textbf{绿盟科技}——安全试验设施与大规模流量分析关键技术。
  \item 相关技术累计服务于电信运营商、国防军队、教育网与公有云，\textbf{惠及数亿终端用户}。
\end{itemize}

% --------------------------------------------------------------- 专利
\section{发明专利}
\note{形成两个专利组合——\textbf{网络测量与安全}（27 项）与\textbf{网络智能}（19 项），
其中绝大多数为第一发明人。代表性专利：}
\begin{itemize}
  \item 基于可编程交换机的高速加密流量双平面特征提取技术。
  \item 基于频域增强与单类自适应学习的加密流量异常检测技术。
  \item 基于内核态与用户态双路径协议特征解析的分布式请求追踪方法。
  \item 基于本地异常感知的元数据驱动型分布式请求追踪方法。
  \item 面向 SD-WAN 的基于 AE-DQN 的双阶段自适应异常流量检测系统。
  \item 深度学习恶意流量检测的预训练数据源审计方法。
  \item 基于 ACME 协议的安全验证方法、系统及相关装置。
  \item 面向跨域分段路由的可信验证方法技术。
  \item 基于带内网络遥测的 OSPF 路由异常检测方法技术。
  \item 一种面向网络快速重路由的完美路由树构建方法。
  \item 互联网测量服务挖掘方法；互联网路由器地理定位方法。
  \item 服务等级目标违规诊断方法、装置、设备及介质。
\end{itemize}

% --------------------------------------------------------------- 学术服务
\section{学术服务}
\begin{itemize}
  \item \textbf{程序委员会共同主席}——ACM SIGCOMM NGNO Workshop 2025。
  \item \textbf{差旅资助共同主席}——ACM APNet 2025。
  \item \textbf{程序委员会委员}——USENIX NSDI 2027（Spring）；USENIX Security 2027；
        IEEE INFOCOM 2023、2024、2025、2027；ACM CoNEXT 2025、2026；
        IEEE ICNP 2024、2025。
  \item \textbf{期刊编委}——《通信技术》。\textbf{客座编辑}——MDPI \emph{Electronics}；
        \emph{IEEE Journal on Selected Areas in Communications}（JSAC）。
  \item \textbf{期刊审稿人}——IEEE/ACM Transactions on Networking；IEEE Transactions on
        Parallel and Distributed Systems；IEEE Transactions on Information Forensics
        and Security。
\end{itemize}

% -------------------------------------------------------------- 教学
\section{教学工作}
\begin{kvlist}
  \kvitem{2023、2024 秋季学期}{\textbf{互联网工程}——本科生课程，32 学时。
    \note{涵盖路由、流量工程、可观测性与互联网测量，含实验环节。}}
  \kvitem{2022--2024 秋季学期}{\textbf{网络优化理论与方法}——研究生课程，32 学时。}
  \kvitem{2022--2025 春季学期}{\textbf{网络异常检测}——研究生课程，32 学时。}
\end{kvlist}

% -------------------------------------------------------------- 学生
\section{指导学生}
\subsection*{在读（13 人）}
\noindent\small
\begin{tabularx}{\linewidth}{@{}>{\raggedright\arraybackslash}p{3.6cm}>{\raggedright\arraybackslash}p{2.7cm}X@{}}
赵铠阳 & 博士·五年级 & SRv6 探测（ICNP'25、ToN'26）；与尹霞教授、吴建平院士联合指导 \\
耿彦涛 & 博士·四年级 & DeepTrace 负责人（SIGCOMM'25、ICSE'26） \\
高杨宸严 & 博士·四年级 & 与赵有健教授联合指导 \\
刘雪枫 & 博士·三年级 & 网络验证 \\
肖泽基 & 博士·三年级 & 与王继龙教授联合指导 \\
Boyang Wang & 博士·二年级 & 与尹霞教授联合指导 \\
童彦图 & 博士·一年级 & SRv6 / 网络测量 \\
李\ 卓 & 硕士·三年级 & 加密流量入侵检测 \\
Hanyang Lin & 硕士·二年级 & eBPF 可观测性 \\
武志衡 & 硕士·二年级 & 分布式追踪 \\
钱林强 & 硕士·一年级 & 可编程数据平面 \\
佟\ 瑶 & 硕士·一年级 & 与尹霞教授联合指导 \\
彭可心 & 硕士·一年级 & 与赵有健教授联合指导 \\
\end{tabularx}

\subsection*{已毕业（6 人）}
\noindent\small
\begin{tabularx}{\linewidth}{@{}>{\raggedright\arraybackslash}p{3.6cm}>{\raggedright\arraybackslash}p{2.7cm}X@{}}
沈俊贤 & 博士 & 华南理工大学 助理教授 \note{（SIGCOMM'23、'25；ToN'26）} \\
Heng Yu & 博士 & 中关村实验室 \note{（TPDS'23）} \\
Yi Qiao & 博士 & 中国航天科技集团 \note{（INFOCOM'24；TPDS'22）} \\
Shangbin Han & 博士 & 泉城实验室 \note{（TIFS'21）} \\
汤之韫 & 硕士 & 华为技术有限公司 \\
李嘉睿 & 硕士 & 阿里云 \\
\end{tabularx}

\vspace{6pt}
\noindent\note{正在招收 2027 年秋季入学的博士 / 硕士研究生，具备系统编程经验者尤佳。}

\end{document}
